\chapter{Teoremas avanzados de límites e inducción}

%----------------------------------------------------------------------
\section{Introducción: Inducción Matemática}
%----------------------------------------------------------------------

\textbf{Simplicio:} Oye, Sócrates, he estado pensando en algo que me ha
estado molestando. Cuando demostramos cosas en matemáticas, normalmente
probamos un enunciado específico. Pero, ¿qué pasa si quiero demostrar que
algo es verdad para todos los números naturales? Por ejemplo, ¿cómo
demuestro un enunciado para $0, 1, 2, 3, 4,$ y así sucesivamente\ldots
para siempre?

\textbf{Sócrates:} ¡Ah, una excelente pregunta! Tienes razón en que no
podemos simplemente comprobar cada caso uno por uno—eso llevaría un tiempo
infinito. Dime, Simplicio, ¿alguna vez has subido una escalera?

\textbf{Simplicio:} ¡Claro que sí! ¿Pero qué tiene eso que ver con esto?

\textbf{Sócrates:} Bueno, imagina una escalera infinitamente alta que
llega hasta el cielo. Si quisiera convencerte de que puedes subir a
cualquier peldaño de esta escalera, ¿qué tendría que mostrarte?

\textbf{Simplicio:} Mmm\ldots supongo que tendrías que demostrarme que
puedo alcanzar el peldaño más bajo, ¿no?

\textbf{Sócrates:} ¡Buen comienzo! ¿Y qué más?

\textbf{Simplicio:} Bueno, si estoy parado en cualquier peldaño
particular, necesitaría saber que puedo alcanzar el siguiente. Así que, si
siempre puedo pasar de un peldaño al siguiente\ldots

\textbf{Sócrates:} ¡Exactamente! Entonces, si puedes alcanzar el primer
peldaño, y siempre puedes pasar del peldaño $k$ al peldaño $k+1$, ¿qué
puedes concluir?

\textbf{Simplicio:} ¡Oh! ¡Entonces puedo alcanzar cualquier peldaño que
quiera! Si quiero llegar al peldaño 100, simplemente empiezo en el peldaño
0, subo al peldaño 1, luego al 2, y sigo subiendo hasta llegar al 100. ¡Y
lo mismo funciona para cualquier número!

\textbf{Sócrates:} ¡Precisamente! Esta es la esencia de la inducción
matemática. Para demostrar que algo es verdad para todos los números
naturales $n$, necesitas exactamente dos cosas:

\begin{itemize}
  \item \textbf{Caso base:} demuestra que es verdad para $n = 0$.
  \item \textbf{Paso inductivo:} demuestra que si es verdad para $n = k$,
        entonces es verdad para $n = k+1$.
\end{itemize}

\textbf{Simplicio:} Espera, pero en el paso inductivo, ¿no estamos
asumiendo lo que queremos probar? ¿No es eso un razonamiento circular?

\textbf{Sócrates:} ¡Una observación perspicaz! Pero no, no es circular.
No estamos asumiendo que el enunciado es verdad para todo $n$. Solo
estamos asumiendo que es verdad para un valor particular $k$, y usando
esa suposición para demostrar que es verdad para $k+1$. Estamos probando
una implicación: ``si $P(k)$ entonces $P(k+1)$''. Combinado con el caso
base, esto crea una reacción en cadena que alcanza cualquier número
natural.

\textbf{Simplicio:} Hmm, creo que lo entiendo. Entonces, ¿la suposición
``es verdad para $k$'' se llama hipótesis inductiva?

\textbf{Sócrates:} ¡Exactamente! Y esa hipótesis es tu herramienta más
poderosa. Es como tener un punto de apoyo en el peldaño $k$ desde el que
puedes impulsarte para llegar al peldaño $k+1$.

\textbf{Simplicio:} Pero, ¿por qué funciona esto? Es decir, ¿por qué
debo creer en este principio?

\textbf{Sócrates:} ¡Ah, una pregunta profunda! Viene de la propia
definición de los números naturales. ¿Cómo crees que se construyen los
números naturales?

\textbf{Simplicio:} Bueno\ldots supongo que comenzamos con 0. Y luego
tenemos 1, que es $0+1$. Y 2 es $1+1$. Así que cada número es el
``sucesor'' del anterior.

\textbf{Sócrates:} ¡Perfecto! Los números naturales están definidos
precisamente por este proceso:

\begin{itemize}
  \item Cero es un número natural.
  \item Si $k$ es un número natural, entonces $k+1$ también lo es.
  \item Estos son los únicos números naturales.
\end{itemize}

¿Ves cómo esto refleja la estructura de la inducción?

\textbf{Simplicio:} ¡Oh! El caso base corresponde a ``cero es un número
natural'', y el paso inductivo corresponde a ``si $k$ es un número
natural, entonces $k+1$ también lo es''. ¡La inducción es simplemente la
construcción de los números naturales convertida en una técnica de
demostración!

\textbf{Sócrates:} ¡Precisamente! No hay ``huecos'' en los números
naturales—ningún número que no pueda alcanzarse partiendo de 0 y sumando
1 repetidamente. Por eso funciona la inducción.

De hecho, esta no es solo una observación filosófica—¡es literalmente
cómo se implementan los números naturales en Lean! En la teoría de tipos
de Lean, un número natural se define inductivamente como:

\begin{itemize}
  \item \mintinline{lean}{zero : ℕ}, el caso base, o
  \item \mintinline{lean}{succ n : ℕ}, el sucesor de otro número natural
        $n$.
\end{itemize}

Así es como se ve esto en el núcleo de Lean:

\begin{minted}{lean}
inductive Nat where
  | zero : Nat
  | succ (n : Nat) : Nat
\end{minted}

Simplemente declaras la existencia de un número natural llamado
\texttt{zero}, y luego declaras que, dado cualquier número natural
$n$, existe otro llamado \texttt{succ n}. (La palabra \texttt{succ} es
solo un nombre; podríamos haberlo llamado \texttt{Alice n}. Lo importante
es que estamos dando una manera de construir un nuevo número natural a
partir de uno ya existente.)

Así, el número 3, por ejemplo, está representado literalmente como
\mintinline{lean}{succ (succ (succ zero))}. ¡El principio de inducción no
solo se parece a esta construcción, sino que lo explota directamente!
Cuando demuestras algo por inducción en Lean, estás trabajando con la
estructura computacional real de cómo los números naturales existen en el
sistema.

\textbf{Simplicio:} ¡Es increíble! ¿Entonces la inducción no es solo
una técnica de demostración, sino que está integrada en el tejido mismo de
cómo Lean entiende los números?

\textbf{Sócrates:} Exactamente. El principio de inducción matemática es
un teorema en muchos marcos matemáticos, pero en la teoría de tipos es una
consecuencia fundamental de cómo están definidos los números naturales.

\textbf{Simplicio:} Bien, estoy convencido de que esta es una técnica de
demostración legítima. ¿Puedes darme un ejemplo?

\textbf{Sócrates:} Por supuesto.

%----------------------------------------------------------------------
\section{Jefe Final: Recíprocos de Sucesiones Convergentes}
%----------------------------------------------------------------------

Uno de los teoremas de límites más importantes concierne a los recíprocos:
si una sucesión converge a un límite no nulo, entonces la sucesión de
recíprocos converge al recíproco del límite. Este resultado es crucial para
demostrar teoremas sobre cocientes y funciones racionales.

Este es un nivel de \emph{Jefe Final}—requiere sintetizar múltiples
técnicas: trabajar con límites no nulos, manipular expresiones algebraicas
complejas y elegir cuidadosamente las estrategias con épsilon.

\subsection{El Marco Matemático}

\begin{theorem}
Sea $a : \mathbb{N} \to \mathbb{R}$ una sucesión que converge a $L$ con
$L \neq 0$, y sea $b : \mathbb{N} \to \mathbb{R}$ definida por
$b(n) = 1/a(n)$ para todo $n$. Entonces $b$ converge a $1/L$.
\end{theorem}

Esta es la demostración técnicamente más desafiante de esta serie de
conferencias.

\subsection{Nuevas Herramientas}

\textbf{abs\_div:} Para cualesquiera números reales $x$ e $y$ (con
$y \neq 0$), se tiene $|x/y| = |x|/|y|$.

\textbf{nonzero\_of\_abs\_pos:} Si $0 < |x|$, entonces $x \neq 0$.

\subsection{Enfoque Estratégico}

Los desafíos clave son:

\begin{itemize}
  \item Asegurar que $a(n) \neq 0$ eventualmente, para que los recíprocos
        estén bien definidos.
  \item Acotar $|1/a(n) - 1/L|$ obteniendo un denominador común.
  \item Elegir el épsilon correcto al aplicar la convergencia de $a$ hacia
        $L$.
  \item Usar la cota inferior sobre $|a(n)|$ para controlar los
        recíprocos.
\end{itemize}

La percepción crítica es que el ``épsilon correcto'' no es la elección
obvia. Usamos $\varepsilon \cdot |L|^2 / 2$, que está diseñado
precisamente para que las desigualdades finales funcionen.

\subsection{Solución en Lean}

\begin{minted}{lean}
Statement InvLim (a : ℕ → ℝ) (L : ℝ) (aToL : SeqLim a L)
    (LneZero : L ≠ 0) (b : ℕ → ℝ)
    (bEqInva : ∀ n, b n = 1 / a n) :
    SeqLim b (1 / L) := by
  choose NhalfL hNhalfL using EventuallyGeHalfLimPos a L
    aToL LneZero
  intro ε hε
  have : 0 < |L| := by apply abs_pos_of_nonzero LneZero
  specialize aToL (ε * |L| * |L| / 2) (by bound)
  choose Na hNa using aToL
  use Na + NhalfL
  intro n hn
  specialize bEqInva n
  rewrite [bEqInva]
  have hnHalfL : NhalfL ≤ n := by bound
  have hna : Na ≤ n := by bound
  specialize hNhalfL n hnHalfL
  specialize hNa n hna
  have : 0 < |a n| := by bound
  have : a n ≠ 0 := by apply nonzero_of_abs_pos this
  have l1 : |1 / a n - 1 / L| = |(L - a n) / (a n * L)| :=
    by field_simp
  have l2 : |(L - a n) / (a n * L)| = |L - a n| / |a n * L| :=
    by apply abs_div
  have l3 : |L - a n| / |a n * L| = |L - a n| / (|a n| * |L|) :=
    by bound
  have l4 : |L - a n| = |-(a n - L)| := by ring_nf
  have l5 : |-(a n - L)| = |a n - L| := by apply abs_neg
  have l6 : |a n - L| / (|a n| * |L|) <
      (ε * |L| * |L| / 2) / (|a n| * |L|) := by
    field_simp; bound
  have l10 : |L - a n| / (|a n| * |L|) =
      |-(a n - L)| / (|a n| * |L|) := by rewrite [l4]; rfl
  have l11 : |-(a n - L)| / (|a n| * |L|) =
      |a n - L| / (|a n| * |L|) := by rewrite [l5]; rfl
  have l13 : ε * |L| * |L| / 2 / (|a n| * |L|) =
      ε * |L| / 2 / |a n| := by field_simp
  have l14 : ε * |L| / 2 / |a n| ≤ ε := by
    field_simp; bound
  linarith [l1, l2, l3, l10, l11, l6, l13, l14]
\end{minted}

\subsection{Demostración en Lenguaje Natural}

\begin{proof}
Supongamos que $a(n) \to L$ con $L \neq 0$, y sea $b(n) = 1/a(n)$.
Debemos demostrar que $b(n) \to 1/L$.

Primero, usando el teorema \texttt{EventuallyGeHalfLimPos}, existe $N_1$
tal que para todo $n \geq N_1$ se tiene $|a(n)| \geq |L|/2 > 0$, lo
que garantiza que $a(n) \neq 0$ y $b(n)$ está bien definida.

Dado $\varepsilon > 0$, como $a(n) \to L$, existe $N_2$ tal que para
todo $n \geq N_2$:
\[
  |a(n) - L| < \frac{\varepsilon |L|^2}{2}.
\]

Sea $N = N_1 + N_2$. Para cualquier $n \geq N$:
\begin{align*}
  |b(n) - 1/L|
  &= \left|\frac{1}{a(n)} - \frac{1}{L}\right| \\
  &= \left|\frac{L - a(n)}{a(n) \cdot L}\right| \\
  &= \frac{|L - a(n)|}{|a(n)| \cdot |L|} \\
  &= \frac{|a(n) - L|}{|a(n)| \cdot |L|} \\
  &< \frac{\varepsilon |L|^2 / 2}{|a(n)| \cdot |L|} \\
  &= \frac{\varepsilon |L|}{2\,|a(n)|} \\
  &\leq \frac{\varepsilon |L|}{2 \cdot |L|/2} \\
  &= \varepsilon.
\end{align*}

Por lo tanto, $b(n) \to 1/L$.
\end{proof}

\subsection{Aplicaciones y Extensiones}

Con este teorema, ahora disponemos de un conjunto completo de herramientas
para los límites de funciones racionales. Combinado con los resultados
sobre sumas y productos, podemos demostrar:

\begin{theorem}[Cociente]
Si $a(n) \to L$ y $c(n) \to M$ con $M \neq 0$, entonces
$a(n)/c(n) \to L/M$.
\end{theorem}

La demostración es directa: primero se muestra que $1/c(n) \to 1/M$
usando el teorema del recíproco, y luego se aplica el teorema del
producto para demostrar que $a(n) \cdot (1/c(n)) \to L \cdot (1/M) = L/M$.

Esto completa la aritmética fundamental de los límites: sumas, productos y
cocientes—los bloques de construcción para analizar polinomios, funciones
racionales y expresiones complejas a lo largo del cálculo y el análisis.

%----------------------------------------------------------------------
\section{Por Casos: Análisis de Casos Sin Hipótesis}
%----------------------------------------------------------------------

Cuando ya tienes una hipótesis \texttt{h : P ∨ Q}, puedes usar
\texttt{cases'} para dividir el objetivo en dos sub-objetivos. Pero,
¿cómo divides en casos cuando no tienes una hipótesis de ese tipo? Ahí es
donde entra la táctica \texttt{by\_cases}.

\subsection{La Táctica \texttt{by\_cases}}

La táctica \texttt{by\_cases} tiene la sintaxis
\texttt{by\_cases h : hecho}, donde \texttt{h} es el nombre para una
nueva hipótesis y \texttt{hecho} es el hecho afirmado en ella. Llamar a
\texttt{by\_cases} crea dos sub-objetivos:

\begin{itemize}
  \item Uno con la hipótesis adicional \texttt{h : hecho}.
  \item Otro con la hipótesis \texttt{h : ¬hecho}.
\end{itemize}

Esto permite realizar un análisis de casos sobre cualquier proposición,
no solo aquellas que ya tienes como hipótesis.

\subsection{Aplicación: Eventualmente Mayor que la Mitad del Límite}

Podemos usar \texttt{by\_cases} para demostrar una versión más general de
\texttt{EventuallyGeHalfLimPos} que no requiere la suposición $L \neq 0$.

\begin{theorem}
Sea $a : \mathbb{N} \to \mathbb{R}$ una sucesión que converge a $L$ (sin
suponer que $L \neq 0$). Entonces existe $N$ tal que para todo
$n \geq N$ se tiene $|a(n)| \geq |L|/2$.
\end{theorem}

\subsection{Solución en Lean}

\begin{minted}{lean}
Statement EventuallyGeHalfLim (a : ℕ → ℝ) (L : ℝ)
    (aToL : SeqLim a L) :
    ∃ N, ∀ n ≥ N, |L| / 2 ≤ |a n| := by
  by_cases h : L = 0
  · use 1
    intro n hn
    rewrite [h]
    bound
  · apply EventuallyGeHalfLimPos a L aToL h
\end{minted}

\subsection{Demostración en Lenguaje Natural}

\begin{proof}
Procedemos por casos según si $L = 0$ o no.

\textbf{Caso 1:} Si $L = 0$, entonces $|L|/2 = 0$. Como $|a(n)| \geq 0$
para todo $n$, podemos tomar $N = 1$ y la desigualdad $|L|/2 \leq |a(n)|$
se cumple trivialmente para todo $n \geq N$.

\textbf{Caso 2:} Si $L \neq 0$, podemos aplicar el teorema
\texttt{EventuallyGeHalfLimPos} (que requiere $L \neq 0$ como hipótesis)
para obtener el $N$ deseado.
\end{proof}

\subsection{Importancia}

La táctica \texttt{by\_cases} es esencial para manejar proposiciones cuyo
valor de verdad importa pero no se conoce de antemano. Nos permite:

\begin{itemize}
  \item Manejar casos límite de manera sistemática.
  \item Demostrar teoremas con hipótesis más débiles.
  \item Aplicar distintas estrategias de demostración según el caso.
\end{itemize}

Esta técnica es omnipresente en el razonamiento matemático, desde el
análisis hasta el álgebra y la combinatoria.

%----------------------------------------------------------------------
\section{Inducción Matemática}
%----------------------------------------------------------------------

La inducción es una de las técnicas de demostración más poderosas para
enunciados sobre números naturales. Permite demostrar infinitos enunciados
(uno para cada número natural) usando solo dos pasos: un caso base y un
paso inductivo.

\subsection{El Principio de Inducción Matemática}

Para demostrar que una propiedad $P(n)$ se cumple para todos los números
naturales $n$, basta con demostrar:

\begin{enumerate}
  \item \textbf{Caso Base:} $P(0)$ es verdadero.
  \item \textbf{Paso Inductivo:} Para todo $k$, si $P(k)$ es verdadero,
        entonces $P(k+1)$ es verdadero.
\end{enumerate}

La intuición es como subir una escalera infinita: si puedes alcanzar el
primer peldaño (caso base), y si estar en cualquier peldaño te permite
alcanzar el siguiente (paso inductivo), entonces puedes alcanzar todos los
peldaños.

\subsection{La Táctica \texttt{induction'}}

La sintaxis de la inducción en Lean es:
\texttt{induction' n with k hk}. Esto significa:

\begin{itemize}
  \item Aplicar inducción sobre la variable $n$.
  \item Usar $k$ como nueva variable muda en el paso inductivo.
  \item Usar \texttt{hk} como hipótesis de inducción sobre $k$.
\end{itemize}

\subsection{Nueva Herramienta: \texttt{ge\_one\_of\_nonzero}}

Si un número natural $n \neq 0$, entonces $1 \leq n$. Este lema sencillo
es útil para manejar el caso en que hemos superado el caso base.

\subsection{Ejemplo: Crecimiento Exponencial}

\begin{theorem}
Para todo número natural $n$, se tiene $n < 2^n$.
\end{theorem}

Este teorema captura la idea de que las funciones exponenciales crecen
más rápido que las lineales—un hecho fundamental con aplicaciones en toda
la matemática y la informática.

\subsection{Solución en Lean}

\begin{minted}{lean}
Statement (n : ℕ) : n < 2 ^ n := by
  induction' n with k hk
  · norm_num
  · by_cases hk0 : k = 0
    · rewrite [hk0]
      norm_num
    · have : 1 ≤ k := by apply ge_one_of_nonzero hk0
      have f1 : k + 1 ≤ 2 * k := by bound
      have f2 : 2 * k < 2 * 2 ^ k := by linarith [hk]
      have f3 : 2 * 2 ^ k = 2 ^ (k + 1) := by ring_nf
      linarith [f1, f2, f3]
\end{minted}

\subsection{Demostración en Lenguaje Natural}

\begin{proof}
Procedemos por inducción sobre $n$.

\textbf{Caso Base} ($n = 0$): Tenemos $0 < 2^0 = 1$, lo cual es
verdadero.

\textbf{Paso Inductivo:} Supongamos que $k < 2^k$ (hipótesis de
inducción). Debemos demostrar que $k + 1 < 2^{k+1}$.

Procedemos por casos según si $k = 0$.

Si $k = 0$, entonces $k + 1 = 1 < 2 = 2^1 = 2^{k+1}$.

Si $k \neq 0$, entonces $k \geq 1$. Por lo tanto:
\begin{align*}
  k + 1 &\leq 2k   && \text{(pues } k \geq 1\text{)} \\
        &< 2 \cdot 2^k && \text{(por hipótesis de inducción)} \\
        &= 2^{k+1}.
\end{align*}

Esto completa la inducción.
\end{proof}

\subsection{Aplicaciones}

La inducción matemática es esencial para demostrar:

\begin{itemize}
  \item Fórmulas para sumas: $\displaystyle\sum_{i=1}^{n} i =
        \frac{n(n+1)}{2}$.
  \item Desigualdades: $n! > 2^n$ para $n$ suficientemente grande.
  \item Propiedades de divisibilidad: $n^3 - n$ es divisible por 6.
  \item Definiciones recursivas: números de Fibonacci, factoriales, etc.
  \item Corrección de algoritmos: demostrar invariantes de bucles.
\end{itemize}

La combinación de la inducción con otras técnicas de demostración (como el
análisis de casos) crea un poderoso conjunto de herramientas para las
matemáticas discretas y la informática.

\subsection{Idea Clave}

Esta demostración ilustra una estrategia importante: dentro de una
demostración inductiva, puede ser necesario realizar análisis de casos
adicionales (usando \texttt{by\_cases}). La hipótesis de inducción te da
ventaja, pero a menudo hay que ser ingenioso en cómo aplicarla,
especialmente cuando se tratan casos frontera o cuando el álgebra requiere
restricciones adicionales.

%%% Local Variables:
%%% mode: LaTeX
%%% TeX-master: "Analisis_en_Lean"
%%% End:
